\begin{table}[!htbp]
\centering
\small
\caption{SEER-HCC income analysis: adjusted stratified Cox model}
\begin{tabular}{llll}
\toprule
Variable & HR & 95\% CI & P-value \\
\midrule
Income Group (High) & 0.765 & (0.690, 0.849) & \textless{}0.001 \\
Age & 1.013 & (1.008, 1.018) & \textless{}0.001 \\
Race (Black) & 0.886 & (0.762, 1.029) & 0.112 \\
Race (Asian or Pacific Islander) & 0.844 & (0.725, 0.983) & 0.029 \\
Sex (Male) & 1.089 & (0.968, 1.226) & 0.157 \\
N Stage (N1) & 1.336 & (1.124, 1.588) & 0.001 \\
M Stage (M1) & 1.997 & (1.705, 2.338) & \textless{}0.001 \\
Tumor Size & 1.004 & (1.003, 1.005) & \textless{}0.001 \\
\bottomrule
\end{tabular}
\begin{minipage}{0.95\linewidth}
\vspace{0.5em}
\footnotesize N = 4668; events = 1608. HR: adjusted hazard ratio; CI: confidence interval. P-values are Wald tests. Categorical reference levels: Income Group = Low; Race = White; Sex = Female; N Stage = N0; M Stage = M0. Age and tumor size HRs are per year and per mm, respectively. Separate baseline hazards are used for T stage, tumor grade, and AFP status; these stratification variables have no estimated HRs.
\end{minipage}
\end{table}
